% Auto-generated from book/00_preview.md — do not edit by hand
% Regenerated by scripts/md_to_latex.py

\chapter{Chapter 0 --- A Preview}
\label{ch:ch00_preview}

This chapter is a guided tour, not a full development. By the end you should see a single chain of ideas:

\[
\text{Pythagorean triples}
\;\longrightarrow\;
\text{sums of two squares}
\;\longrightarrow\;
\text{four squares / quaternions}
\;\longrightarrow\;
S^3
\;\longrightarrow\;
\text{Hopf fibration}
\;\longrightarrow\;
\text{gauged lattice \& flux flywheels}
\;\longrightarrow\;
\text{\(Z\mapsto\) flux map (model)}
\]

and know which links are classical mathematics and which are working models or hypotheses from Kingdom Come.

\textbf{Learning goals}

\begin{enumerate}
\item Recall why geometry already sits inside elementary number theory.  
\item Meet unit quaternions as \(S^3\) and the Hopf map \(S^3\to S^2\).  
\item Glimpse flux flywheels and the \(Z\)-to-flux idea without full machinery.  
\item Spot the \(350/\pi\) signature as a \textbf{Hypothesis}, not a theorem.  
\item Own a visual roadmap for the rest of the book.
\end{enumerate}

\textbf{Figures in this chapter} (all from Kingdom Come assets; paths relative to this file)

\begin{center}
\small
\setlength{\tabcolsep}{4pt}
\begin{tabularx}{\textwidth}{@{}>{\raggedright\arraybackslash}p{0.12\textwidth}>{\raggedright\arraybackslash}X>{\raggedright\arraybackslash}X@{}}
\toprule
Tag & File & Role \\
\midrule
Fig. 0.1 & \texttt{figures/\allowbreak{}fig0\_\allowbreak{}1\_\allowbreak{}hopf\_\allowbreak{}linked\_\allowbreak{}fibers.\allowbreak{}png} & Linked Hopf fibers in stereographic view \\
Fig. 0.2 & \texttt{figures/\allowbreak{}fig0\_\allowbreak{}2\_\allowbreak{}hopf\_\allowbreak{}stereographic.\allowbreak{}png} & Stereographic preview of the fibration \\
Fig. 0.3 & \texttt{figures/\allowbreak{}fig0\_\allowbreak{}3\_\allowbreak{}hopf\_\allowbreak{}bundle.\allowbreak{}png} & Bundle picture: total space \(\to\) base \\
Fig. 0.4 & \texttt{figures/\allowbreak{}fig0\_\allowbreak{}4\_\allowbreak{}flux\_\allowbreak{}flywheel\_\allowbreak{}scales.\allowbreak{}jpg} & Flux flywheel scales / nested structure \\
Fig. 0.5 & \texttt{figures/\allowbreak{}fig0\_\allowbreak{}5\_\allowbreak{}hopf\_\allowbreak{}lattice\_\allowbreak{}pi\_\allowbreak{}cycle.\allowbreak{}jpg} & Gauged lattice meets \(350/\pi\) cycle motif \\
\bottomrule
\end{tabularx}
\end{center}


\begin{center}
\small
\setlength{\tabcolsep}{4pt}
\begin{tabularx}{\textwidth}{@{}>{\raggedright\arraybackslash}p{0.12\textwidth}>{\raggedright\arraybackslash}X>{\raggedright\arraybackslash}X@{}}
\toprule
Aux A0.1 & \texttt{figures/\allowbreak{}aux\_\allowbreak{}z\_\allowbreak{}map\_\allowbreak{}he.\allowbreak{}png} & Helium \(Z=2\) flux-flywheel still (closed-shell baseline) \\
\midrule
Aux A0.2 & \texttt{figures/\allowbreak{}aux\_\allowbreak{}z\_\allowbreak{}map\_\allowbreak{}fe.\allowbreak{}png} & Iron \(Z=26\) flux-flywheel still (mid-table contrast) \\
\bottomrule
\end{tabularx}
\end{center}


\bigskip
\noindent\rule{\textwidth}{0.4pt}
\bigskip

\section{0.1 From Pythagorean triples to sums of two squares}
\label{ch:ch00_preview:0-1-from-pythagorean-triples-to-sums-of-two-squares}

A \textbf{Pythagorean triple} is a triple of positive integers \((a,b,c)\) with

\[
a^2 + b^2 = c^2.
\]

The triple \((3,4,5)\) is the smallest; infinitely many others exist. Already geometry and arithmetic talk: the equation is the integer-point condition on a circle of radius \(c\), or the norm condition \(N(a+bi)=c^2\) in the Gaussian integers \(\mathbb{Z}[i]\).

Hatcher's \emph{Topology of Numbers} opens by pushing this geometric instinct further. Rational slopes, mediants

\[
\frac{a}{c}\oplus\frac{b}{d} \;=\; \frac{a+b}{c+d},
\]

and the \textbf{Farey diagram} organize which fractions ``touch'' which others. Continued fractions become zigzag paths on that diagram. Binary quadratic forms

\[
f(x,y) = ax^2 + bxy + cy^2
\]

acquire \textbf{topographs}---value landscapes on the dual of the Farey triangulation---whose periodic separator lines encode Pell-type infinitude and classification.

We will not rebuild that plane geometry here (read Hatcher Chapters 0--4 for the full story). We only need the moral:

> Elementary number theory is already a theory of \textbf{pictures}: adjacency, paths, periods, and symmetries of integer data.

The present book asks what happens when those pictures are lifted from two variables and the rational/hyperbolic plane to \textbf{four real coordinates} organized by quaternion multiplication and a celebrated fiber bundle.

\bigskip
\noindent\rule{\textwidth}{0.4pt}
\bigskip

\section{0.2 Four squares and the birth of quaternions}
\label{ch:ch00_preview:0-2-four-squares-and-the-birth-of-quaternions}

Lagrange's four-square theorem asserts that every natural number is a sum of four integer squares:

\[
n = a^2 + b^2 + c^2 + d^2.
\]

There is a classical algebraic home for this statement: the \textbf{quaternions}

\[
\mathbb{H} = \{ \, q = a + bi + cj + dk \mid a,b,c,d\in\mathbb{R} \, \},
\]

with

\[
i^2=j^2=k^2=ijk=-1.
\]

The multiplicative \textbf{norm}

\[
N(q) = a^2+b^2+c^2+d^2 = q\,\overline{q}
\]

satisfies \(N(q_1 q_2)=N(q_1)N(q_2)\). Integer quaternions (Lipschitz \(\mathbb{Z}[i,j,k]\), better still the \textbf{Hurwitz order}) turn the four-square theorem into an arithmetic fact about norms---just as Gaussian integers turn two-square questions into unique factorization and primes congruent to \(1\bmod 4\).

\textbf{Claim type:} four-square theorem and the norm identity are \textbf{Theorems} (classical). Their use as the \emph{algebraic backbone} of Kingdom Come's geometry is a \textbf{Model} choice we make explicit.

Chapter 1 develops multiplication, conjugation, unit quaternions, and the Lipschitz/Hurwitz orders in full. For the preview, one geometric fact is enough:

\[
\{ q\in\mathbb{H} : N(q)=1 \} \;=\; S^3 \subset \mathbb{R}^4.
\]

The unit sphere in four dimensions is not an ornament. It is the total space of the map we study next.

\bigskip
\noindent\rule{\textwidth}{0.4pt}
\bigskip

\section{0.3 Circles and spheres as number-theoretic objects}
\label{ch:ch00_preview:0-3-circles-and-spheres-as-number-theoretic-objects}

In Hatcher's world, the circle \(S^1\) already appears indirectly: slopes, rotations in the plane of a binary form, and periodic topograph orbits. The two-sphere \(S^2\) is less prominent in elementary texts, but it is the natural \textbf{base} once we pass to unit quaternions.

Write a unit quaternion as a pair of complex numbers (up to the usual identification)

\[
q \;\longleftrightarrow\; (z_1,z_2)\in\mathbb{C}^2,\qquad |z_1|^2+|z_2|^2=1.
\]

The classical \textbf{Hopf map} sends \(q\) to a point on \(S^2\). One real-coordinate form used in Kingdom Come's documentation is

\begin{align*}
y_1 &= x_1^2 - x_2^2,\\
y_2 &= 2 x_1 x_2,\\
y_3 &= 2(x_3 x_4 + x_1 x_2),
\end{align*}

for a unit 4-vector \((x_1,x_2,x_3,x_4)\in S^3\) (see \texttt{kingdom} theory notes and \texttt{flux\_\allowbreak{}hopf\_\allowbreak{}lib}). Equivalent complex forms identify the base point with a ratio \([z_1:z_2]\in\mathbb{CP}^1\cong S^2\).

\textbf{Fiber intuition.} Over each base point sits a \textbf{circle's worth} of unit quaternions---the fiber---obtained by rotating a phase along that circle. Distinct fibers are \textbf{linked}: any two are linked once (Hopf invariant \(1\)). Stereographic projection \(S^3\setminus\{\mathrm{pt}\}\to\mathbb{R}^3\) displays those fibers as linked Villarceau circles on nested tori.

\begin{figure}[htbp]
  \centering
  \includegraphics[width=0.92\textwidth,height=0.42\textheight,keepaspectratio]{fig0_1_hopf_linked_fibers.png}
  \caption{Figure 0.1 --- Linked Hopf fibers (stereographic view). Source: Kingdom Come \texttt{app/\allowbreak{}assets/\allowbreak{}home/\allowbreak{}hopf\_\allowbreak{}linked\_\allowbreak{}fibers.\allowbreak{}png}.}
  \label{fig:fig0-1-hopf-linked-fibers}
\end{figure}

\begin{quote}\small\textit{Figure 0.1.* Two (or more) fibers of the Hopf fibration after stereographic projection to \(\mathbb{R}^3\). Linking is the topological signature that will later underwrite \textbf{protection} of flux configurations: you cannot unlink fibers by continuous deformation without cutting.}\end{quote}

\begin{figure}[htbp]
  \centering
  \includegraphics[width=0.92\textwidth,height=0.42\textheight,keepaspectratio]{fig0_2_hopf_stereographic.png}
  \caption{Figure 0.2 --- Stereographic Hopf preview. Source: Kingdom Come \texttt{app/\allowbreak{}assets/\allowbreak{}hopf\_\allowbreak{}preview.\allowbreak{}png}.}
  \label{fig:fig0-2-hopf-stereographic}
\end{figure}

\begin{quote}\small\textit{Figure 0.2.* A portal-style stereographic preview of the fibration: the same total space seen as a single composed picture rather than isolated fibers. In the Gradio app this is the ``Hopf Visualizer'' family of views (2D projections are preferred on CPU-only hosts).}\end{quote}

Hatcher's Farey diagram arranges \textbf{rational slopes} so that adjacency and mediants are visible. Hopf arranges \textbf{circle fibers} over a sphere so that linking and phase are visible. The analogy is not identity of objects; it is identity of \textbf{method}: a picture that makes the right discrete or continuous relationships obvious.

\textbf{Forward pointer.} Calling Hopf a ``Farey analogue'' is a \textbf{Model} metaphor only. \textbf{Chapter 3 will turn this metaphor into precise discrete definitions}---adjacency, mediants, and gauge data on a lattice in \(S^3\)---so that the analogy becomes a construction rather than a slogan.

\bigskip
\noindent\rule{\textwidth}{0.4pt}
\bigskip

\section{0.4 The Hopf fibration as a higher-dimensional Farey analogue}
\label{ch:ch00_preview:0-4-the-hopf-fibration-as-a-higher-dimensional-farey-analogu}

We state the comparison carefully so it does not pretend to be a theorem of uniqueness.

\begin{center}
\small
\setlength{\tabcolsep}{4pt}
\begin{tabularx}{\textwidth}{@{}>{\raggedright\arraybackslash}X>{\raggedright\arraybackslash}X@{}}
\toprule
Farey diagram (Hatcher) & Hopf geometry (this book) \\
\midrule
Vertices: rationals / cusps & Points of a discrete set in \(S^3\) (or in \(\mathbb{H}\)) \\
Edges: adjacency of fractions & Fiber segments / lattice bonds projecting under Hopf \\
Mediant of two fractions & Quaternionic (or lattice) combination rule, \textbf{to be axiomatized in Ch. 3} \\
Paths: continued fractions & Phase winding along fibers; walks on the gauged lattice \\
Symmetries: linear fractional maps & Left/right multiplication by unit quaternions; gauge actions \\
\bottomrule
\end{tabularx}
\end{center}


\begin{figure}[htbp]
  \centering
  \includegraphics[width=0.92\textwidth,height=0.42\textheight,keepaspectratio]{fig0_3_hopf_bundle.png}
  \caption{Figure 0.3 --- Hopf fibration as a bundle. Source: Kingdom Come \texttt{app/\allowbreak{}assets/\allowbreak{}home/\allowbreak{}hopf\_\allowbreak{}fibration\_\allowbreak{}bundle.\allowbreak{}png}.}
  \label{fig:fig0-3-hopf-bundle}
\end{figure}

\begin{quote}\small\textit{Figure 0.3.* Conceptual bundle picture: total space \(S^3\), base \(S^2\), fibers circles. In Part II we \textbf{discretize} this picture---placing a lattice in the total space, projecting to the base, and defining adjacency---so that Farey-like combinatorial operations become available again.}\end{quote}

\textbf{Claim type:} Hopf fibration, linking, and the double cover \(\mathrm{Spin}(3)\to SO(3)\) via unit quaternions are \textbf{Theorems} of classical topology/geometry. Calling Hopf a ``Farey analogue'' remains a \textbf{Model} metaphor until Chapter 3 supplies the discrete axioms (see also Open Problem 1: \emph{canonical quaternionic Farey structure}).

\bigskip
\noindent\rule{\textwidth}{0.4pt}
\bigskip

\section{0.5 First sight of the gauged Hopf lattice and flux flywheels}
\label{ch:ch00_preview:0-5-first-sight-of-the-gauged-hopf-lattice-and-flux-flywheel}

Kingdom Come's physical \textbf{Model} (not yet a theorem of nature) can be stated in one paragraph:

> Physics emerges from \textbf{topologically protected flux structures} on a \textbf{gauged Hopf lattice} embedded in a \textbf{porous vacuum}. The Hopf fibration supplies the geometric backbone; quaternions supply the algebra; stable rotating configurations---\textbf{flux flywheels}---anchor emergent matter. Observer-linked phase holonomy between fibers damps toward synchrony with a coupling scale \(\kappa\).

A \textbf{flux flywheel}, for preview purposes, is a lattice-supported circulating flux whose topology (linking, winding, charge) obstructs continuous decay to a trivial configuration. In Hatcher's topographs, \textbf{periodicity} of separator lines produces infinite families of representations (Pell). Here, \textbf{periodicity of a flywheel} is the dynamical counterpart: a closed orbit of gauge and phase that refuses to unwind.

\begin{figure}[htbp]
  \centering
  \includegraphics[width=0.92\textwidth,height=0.42\textheight,keepaspectratio]{fig0_4_flux_flywheel_scales.jpg}
  \caption{Figure 0.4 --- Flux flywheel scales. Source: Kingdom Come \texttt{app/\allowbreak{}assets/\allowbreak{}bitcoin\_\allowbreak{}pi/\allowbreak{}flux\_\allowbreak{}flywheel\_\allowbreak{}scales.\allowbreak{}jpg}.}
  \label{fig:fig0-4-flux-flywheel-scales}
\end{figure}

\begin{quote}\small\textit{Figure 0.4.* Nested / multi-scale flywheel imagery from the Kingdom Come observation assets. Read it as a \textbf{visual metaphor plus design sketch} for the lattice construction of Chapters 3--5: concentric scales, rotational identity, and coupling between levels. Formal definitions of lattice sites, gauge fields, and topological charge come later.}\end{quote}

In code, lattice and flywheel experiments live under:

\begin{itemize}
\item \texttt{kingdom.\allowbreak{}core.\allowbreak{}lattice} --- gauged lattice dynamics and stability comparisons  
\item \texttt{kingdom.\allowbreak{}core.\allowbreak{}flux\_\allowbreak{}flywheel} --- \texttt{map\_\allowbreak{}z\_\allowbreak{}to\_\allowbreak{}flywheel}, \texttt{map\_\allowbreak{}z\_\allowbreak{}to\_\allowbreak{}flywheel\_\allowbreak{}extended}  
\item Gradio tabs: \textbf{Lattice Simulator}, \textbf{Flux Flywheel}
\end{itemize}

Chapter 5 will introduce \textbf{flux topographs}: value landscapes for quaternionic (or lattice flux) functionals, the direct conceptual child of Conway topographs.

\bigskip
\noindent\rule{\textwidth}{0.4pt}
\bigskip

\section{0.6 A first \(Z\)-to-flux sketch and the \(350/\pi\) signature}
\label{ch:ch00_preview:0-6-a-first-z-to-flux-sketch-and-the-350-pi-signature}

\subsection{The \(Z\)-map (\textbf{Model})}
\label{ch:ch00_preview:the-z-map-model}

Kingdom Come implements a map from atomic number \(Z\) to a flywheel configuration and a bundle of stability / chemistry-facing metrics:

\[
Z \;\longmapsto\; \bigl(\text{flywheel state},\;\text{stability score},\;\text{shell-like descriptors},\;\ldots\bigr).
\]

Primary entry points: \texttt{map\_\allowbreak{}z\_\allowbreak{}to\_\allowbreak{}flywheel} and \texttt{map\_\allowbreak{}z\_\allowbreak{}to\_\allowbreak{}flywheel\_\allowbreak{}extended} in \texttt{kingdom.\allowbreak{}core.\allowbreak{}flux\_\allowbreak{}flywheel}. Detuning parameters (frequency offset, gauge strength, depth/layers) select \textbf{stability islands} in parameter space---\textbf{Magic Islands}---including a documented ultra-stable lock used as a noble-gas-like template.

The portal's element tour renders each \(Z\) with shell clouds and a flux ring. Two static stills from that tour are kept beside the main figures:

\begin{figure}[htbp]
  \centering
  \includegraphics[width=0.92\textwidth,height=0.42\textheight,keepaspectratio]{aux_z_map_he.png}
  \caption{Auxiliary Figure A0.1 --- Helium (\(Z=2\)) flux-flywheel still: closed-shell archetype with high model stability. Source: Kingdom Come \texttt{z\_\allowbreak{}knowns/\allowbreak{}frame\_\allowbreak{}0002.\allowbreak{}png}.}
  \label{fig:aux-z-map-he}
\end{figure}

\begin{quote}\small\textit{Auxiliary Figure A0.1 --- Helium (\(Z=2\)).* Flux-flywheel element card from the portal tour: noble-gas closed shell, high stability score in the working model, and a simple shell/flux-ring visualization. Use this as the ``quiet baseline'' when comparing later \(Z\) values. Source: \texttt{z\_\allowbreak{}knowns/\allowbreak{}frame\_\allowbreak{}0002.\allowbreak{}png}.}\end{quote}

\begin{figure}[htbp]
  \centering
  \includegraphics[width=0.92\textwidth,height=0.42\textheight,keepaspectratio]{aux_z_map_fe.png}
  \caption{Auxiliary Figure A0.2 --- Iron (\(Z=26\)) flux-flywheel still: mid-table nuclear/chemical specialness meets a richer flywheel profile. Source: Kingdom Come \texttt{z\_\allowbreak{}knowns/\allowbreak{}frame\_\allowbreak{}0026.\allowbreak{}png}.}
  \label{fig:aux-z-map-fe}
\end{figure}

\begin{quote}\small\textit{Auxiliary Figure A0.2 --- Iron (\(Z=26\)).* Same tour format at iron: a mid-table configuration where real-world nuclear and chemical specialness can later be checked against model metrics (stability, ionization-energy anchors, magnetic-moment proxies in Chapters 7 and 10). Source: \texttt{z\_\allowbreak{}knowns/\allowbreak{}frame\_\allowbreak{}0026.\allowbreak{}png}.}\end{quote}

\textbf{Claim type:} the existence of a coded map \(Z\mapsto\) configuration is a \textbf{Software fact}. Interpreting it as \emph{the} emergence mechanism of the periodic table is a \textbf{Model}. Quantitative alignment with experiment is an empirical research program (Chapters 7 and 10), not a free theorem.

\subsection{The \(350/\pi\) signature (\textbf{Hypothesis})}
\label{ch:ch00_preview:the-350-pi-signature-hypothesis}

Kingdom Come's constants module records a topological clock scale

\[
W_g \;=\; \frac{350}{\pi} \;\approx\; 111.408
\]

(\texttt{WG\_\allowbreak{}FROM\_\allowbreak{}350\_\allowbreak{}OVER\_\allowbreak{}PI}, with a documented lock value \(W_{g,\mathrm{lock}}\) in \texttt{flux\_\allowbreak{}hopf\_\allowbreak{}lib} / \texttt{kingdom.\allowbreak{}core.\allowbreak{}constants}). The same numerical motif appears in several \textbf{observation tracks} of the portal: Bitcoin Pi Cycle narratives, TLS tree branch bursts, cuprate superconductor sketches, pulsar-timing numerics, and related notes.

\begin{figure}[htbp]
  \centering
  \includegraphics[width=0.92\textwidth,height=0.42\textheight,keepaspectratio]{fig0_5_hopf_lattice_pi_cycle.jpg}
  \caption{Figure 0.5 --- Hopf lattice and Pi-cycle motif. Source: Kingdom Come \texttt{app/\allowbreak{}assets/\allowbreak{}bitcoin\_\allowbreak{}pi/\allowbreak{}hopf\_\allowbreak{}lattice\_\allowbreak{}pi\_\allowbreak{}cycle.\allowbreak{}jpg}.}
  \label{fig:fig0-5-hopf-lattice-pi-cycle}
\end{figure}

\begin{quote}\small\textit{Figure 0.5.* Asset bridging the \textbf{gauged Hopf lattice} picture with the \(\pi\)-cycle / \(350/\pi\) observational storyline. In this book it marks the boundary between geometry (Parts I--IV) and the hypothesis layer (Part V).}\end{quote}

\textbf{Claim type: Hypothesis.} Recurrence of \(350/\pi\) across domains is \textbf{not} proved from first principles in this chapter. Chapter 10 will:

\begin{enumerate}
\item list each domain with source and method,  
\item require a statistical validation checklist,  
\item state what would \textbf{confirm, refine, or falsify} the hypothesis.
\end{enumerate}

Until then, treat \(W_g=350/\pi\) as a \textbf{named constant of the working model}, not as a theorem of number theory.

\bigskip
\noindent\rule{\textwidth}{0.4pt}
\bigskip

\section{0.7 Visual roadmap: Hatcher diagrams meet Hopf and Gradio}
\label{ch:ch00_preview:0-7-visual-roadmap-hatcher-diagrams-meet-hopf-and-gradio}

Think of the book as five floors of the same building:

\begin{lstlisting}[style=qga]
Part V   Ch. 10     Observations & validation     <- Hypothesis & Model
Part IV  Ch. 8-9    Class groups & algebras       <- Hatcher Ch. 7-8 lift
Part III Ch. 5-7    Forms, topographs, Z-map      <- Hatcher Ch. 4-6 lift
Part II  Ch. 3-4    Gauged lattice (Farey lift)   <- Hatcher Ch. 1-3 lift
Part I   Ch. 1-2    Quaternions & Hopf            <- foundations (new floor)
         Ch. 0      This preview
\end{lstlisting}

\textbf{How the five main figures will reappear}

\begin{center}
\small
\setlength{\tabcolsep}{4pt}
\begin{tabularx}{\textwidth}{@{}>{\raggedright\arraybackslash}X>{\raggedright\arraybackslash}X@{}}
\toprule
Figure & Returns in \\
\midrule
0.1 Linked fibers & Ch. 2 (definition, Hopf invariant); Ch. 4 (symmetry orbits) \\
0.2 Stereographic preview & Ch. 1--2 (coordinates, charts); Gradio labs \\
0.3 Bundle diagram & Ch. 2--3 (discrete total space / base) \\
0.4 Flywheel scales & Ch. 3, 5, 7 (lattice, topographs, \(Z\)-map) \\
0.5 Lattice / \(\pi\)-cycle & Ch. 3 (lattice); Ch. 10 (observations) \\
\bottomrule
\end{tabularx}
\end{center}


**Portal \(\leftrightarrow\) chapter cheat sheet**

\begin{center}
\small
\setlength{\tabcolsep}{4pt}
\begin{tabularx}{\textwidth}{@{}>{\raggedright\arraybackslash}X>{\raggedright\arraybackslash}X@{}}
\toprule
Gradio / module & First serious chapter \\
\midrule
Hopf Visualizer \(\cdot\) \texttt{core.\allowbreak{}hopf} \(\cdot\) \texttt{viz.\allowbreak{}hopf\_\allowbreak{}plotly} & 2 \\
Quaternion helpers \(\cdot\) \texttt{core.\allowbreak{}quaternion} & 1 \\
Lattice Simulator \(\cdot\) \texttt{core.\allowbreak{}lattice} & 3--4 \\
Flux Flywheel slider \(\cdot\) \texttt{core.\allowbreak{}flux\_\allowbreak{}flywheel} & 5--7 \\
Magic Island \(\cdot\) \texttt{viz.\allowbreak{}magic\_\allowbreak{}island} & 6 \\
Observations tabs \(\cdot\) \texttt{constants.\allowbreak{}WG\_\allowbreak{}FROM\_\allowbreak{}350\_\allowbreak{}OVER\_\allowbreak{}PI} & 10 \\
\bottomrule
\end{tabularx}
\end{center}


\bigskip
\noindent\rule{\textwidth}{0.4pt}
\bigskip

\section{0.8 Reading plan for the rest of the book}
\label{ch:ch00_preview:0-8-reading-plan-for-the-rest-of-the-book}

\begin{enumerate}
\item \textbf{Part I (Ch. 1--2)} makes quaternions algebraic and geometric, then defines the Hopf fibration carefully and rebuilds Figures 0.1--0.3 from first principles.  
\item \textbf{Part II (Ch. 3--4)} constructs the gauged Hopf lattice and its symmetries---the Farey lift (Open Problem 1).  
\item \textbf{Part III (Ch. 5--7)} develops forms, flux topographs, Magic Islands, and the \(Z\)-map as representation theory (Open Problems 2--4).  
\item \textbf{Part IV (Ch. 8--9)} treats composition, class groups, and quaternion algebras (Open Problem 6).  
\item \textbf{Part V (Ch. 10)} collects multi-domain observations (\(350/\pi\), \(Z\mapsto\) correlations, Magic Islands), states validation protocols (Table T4), and summarizes Open Problems OP1--OP6 (Open Problem 5).
\end{enumerate}

If you are impatient for physics, read 0 \(\rightarrow\) 2 \(\rightarrow\) 3 (flywheel sections) \(\rightarrow\) 7 \(\rightarrow\) 10, then return to fill the arithmetic spine. If you are a number theorist first, read 0--9 with Hatcher open, and treat Chapter 10 as a speculative appendix until the checklists are green.

\bigskip
\noindent\rule{\textwidth}{0.4pt}
\bigskip

\section{Exercises (preview level)}
\label{ch:ch00_preview:exercises-preview-level}

\textbf{0.A.} Verify \(N(q_1 q_2)=N(q_1)N(q_2)\) for two random integer quaternions by hand or with \texttt{kingdom.\allowbreak{}core.\allowbreak{}quaternion}.

\textbf{0.B.} In the Gradio Hopf Visualizer, load \textbf{Classic Hopf} and identify which panel corresponds most closely to Figure 0.1 versus Figure 0.2.

\textbf{0.C.} Run \texttt{map\_\allowbreak{}z\_\allowbreak{}to\_\allowbreak{}flywheel} for \(Z\in\{2,10,26,79\}\). Record stability scores; do \textbf{not} yet interpret them as chemical truth---only as model outputs.

\textbf{0.D.} Write one paragraph distinguishing \textbf{Theorem / Model / Hypothesis} for: (i) four-square theorem, (ii) Hopf linking, (iii) \(W_g=350/\pi\) as a multi-domain constant.

\textbf{0.E.} Open Hatcher TN Chapter 0 and list three pictures you expect to have quaternionic or Hopf analogues by the end of our Chapter 5.

\bigskip
\noindent\rule{\textwidth}{0.4pt}
\bigskip

\section{Code and asset pointers}
\label{ch:ch00_preview:code-and-asset-pointers}

\begin{lstlisting}[style=qga]
kingdom.core.quaternion      # Quaternion, from_hopf_coords, hopf_image
kingdom.core.hopf            # sample_fiber, sample_fiber_family
kingdom.core.lattice         # gauged lattice simulations
kingdom.core.flux_flywheel   # map_z_to_flywheel[_extended]
kingdom.core.constants       # WG_FROM_350_OVER_PI, W_G_LOCK, ? defaults
kingdom.viz.hopf_plotly      # portal-style multi-panel figures
\end{lstlisting}

\textbf{Figure sources (Kingdom Come):} see \texttt{figures/\allowbreak{}ATTRIBUTION.\allowbreak{}md}.

\textbf{Parallel reading:} Hatcher, \emph{Topology of Numbers}, Chapter 0; project file \texttt{HATCHER\_\allowbreak{}MAP.\allowbreak{}md}.

\bigskip
\noindent\rule{\textwidth}{0.4pt}
\bigskip


